\section{Introduction}

Vehicle-to-Everything (V2X) communication enables safety-critical and compute-intensive applications such as cooperative perception, autonomous driving, and real-time hazard avoidance. These applications impose heterogeneous Quality-of-Service (QoS) requirements: safety messages may become unusable within approximately 100 ms, whereas compute and infotainment workloads can tolerate substantially larger delays. At the same time, processing the growing volume of camera, radar, and LiDAR data entirely on resource-constrained vehicle Electronic Control Units (ECUs) is increasingly difficult. Vehicular Edge Computing (VEC) addresses this limitation by offloading computation to roadside Multi-access Edge Computing (MEC) servers.

The effectiveness of offloading, however, depends jointly on where computation is performed and how the offloaded workload reaches the edge. No single Radio Access Technology (RAT) is uniformly optimal across task types and mobility conditions. LTE, ITS-G5, WiFi-6, and 5G-NR exhibit different latency, throughput, coverage, and mobility characteristics. A RAT suitable for a delay-tolerant compute task may therefore be unsuitable for a safety-critical task generated by the same vehicle. Moreover, binary local/offload decisions cannot exploit parallel local and MEC processing. These observations motivate treating RAT selection and partial computation offloading as a joint decision rather than as independent optimization stages.

A second challenge emerges as vehicle density increases. Vehicles associated with the same RSU compete not only for wireless resources but also for a shared MEC processing queue. Consequently, aggressive offloading that performs well for an isolated vehicle may degrade when many vehicles simultaneously submit workloads. This degradation can also be hidden by conventional communication metrics. Packet Reception Ratio (PRR) counts a packet as successfully delivered even when it arrives after its application deadline; for safety-critical traffic, however, a late result may be operationally equivalent to a lost one. Evaluation must therefore distinguish successful reception from deadline-compliant delivery.

To address these coupled challenges, we propose MA-ROPE, a multi-agent framework that jointly learns RAT selection and partial computation offloading under shared MEC contention. Each vehicle selects among LTE, ITS-G5, WiFi-6, and 5G-NR while adapting its offloaded workload fraction to channel conditions, task requirements, energy consumption, and queue state. A Dueling Double DQN learns the decision policy, while Federated Averaging (FedAvg) enables cooperative learning through model-parameter exchange without centrally pooling raw vehicle observations.

\textbf{Central research question:} Can a federated multi-RAT offloading policy preserve deadline-compliant V2X delivery as fleet density and shared MEC contention increase?

Our results show that the distinction becomes substantial under dense traffic. Under URLLC stress, MA-ROPE achieves a fleet On-Time delivery Rate (OTR) of 0.834 at $N=32$, compared with 0.525 for single-RAT AllOffload, while consuming 25.2\% less energy. The underlying density experiment shows that MA-ROPE decreases only from 0.887 OTR at $N=8$ to 0.834 at $N=32$, whereas AllOffload decreases from 0.877 to 0.525.

To quantify this density robustness directly, we introduce the Deadline Retention Ratio (DRR), $R_D$, which measures the fraction of reference-density on-time delivery retained as fleet size increases. Across the fourfold increase from $N=8$ to $N=32$, MA-ROPE achieves $R_D=0.940$, whereas single-RAT AllOffload achieves $R_D=0.599$. Thus, MA-ROPE retains approximately 94\% of its reference deadline performance compared with approximately 60\% for single-RAT offloading. The complete definition and density analysis are presented in Section V-F.

Despite extensive research on VEC offloading, existing approaches generally address only subsets of heterogeneous RAT selection, partial computation offloading, multi-vehicle MEC contention, mobility-aware channel behavior, and federated cooperative learning. MA-ROPE brings these dimensions together to evaluate not only performance at a fixed operating point, but also how well deadline-compliant service survives increasing fleet contention.

The main contributions of this work are:

- \textbf{\textit{C1:}} \textbf{Joint multi-RAT partial offloading.} We formulate RAT selection and partial computation offloading as a unified decision across LTE, ITS-G5, WiFi-6, and 5G-NR, incorporating heterogeneous channel behavior and shared RSU queue contention.

- \textbf{\textit{C2:}} \textbf{Federated queue-aware learning.} We develop a Dueling Double DQN policy coordinated through FedAvg, enabling vehicles to learn cooperatively without centrally pooling their raw observations while explicitly accounting for multi-vehicle MEC contention.

- \textbf{\textit{C3:}} \textbf{Deadline-aware density evaluation.} We distinguish communication reliability (PRR) from application-useful on-time delivery (OTR) and introduce the Deadline Retention Ratio $R_D$ to quantify how much deadline performance a policy preserves as fleet density increases. MA-ROPE achieves $R_D=0.940$ over $N=8\rightarrow32$, compared with 0.599 for AllOffload.

- \textbf{\textit{C4:}} \textbf{Dense-fleet performance validation.} Under URLLC stress at $N=32$, MA-ROPE achieves 83.4\% fleet OTR versus 52.5\% for AllOffload, while consuming 25.2\% less energy. Against the centralized MADDPG upper bound, MA-ROPE trades a small $U_{OTR}$ difference for lower energy consumption and decentralized raw-observation handling.
